\documentclass[10pt,a4paper]{book}

\usepackage{amsmath}
\usepackage{amsthm}
\usepackage{amssymb}
\usepackage{amsfonts}
\usepackage{mathtools}
\usepackage{pgfplots}
\usepackage{graphicx}
\usepackage{caption}
\usepackage{float}
\usepackage{multicol}
\usepackage{changepage}
\usepackage{xcolor}
\usepackage{environ}
\usepackage[top=25mm,bottom=26mm]{geometry}
\usepackage{lmodern}
\usepackage[utf8]{inputenc}
\usepackage[T1,T2A]{fontenc}
\usepackage[greek,belarusian]{babel}
\usepackage{fancyhdr}
\usepackage{tcolorbox}
\usepackage{fontspec}

\usepackage{tikz}
\usetikzlibrary{angles,quotes,babel,calc}

\usepackage[explicit]{titlesec}
\setmainfont{Exo2-Light.ttf}
\newfontfamily\captfont{Exo2-Black.ttf}
\newfontfamily\medfont{Exo2-Medium.ttf}
\newfontfamily\figfont{Exo2-BoldItalic.ttf}
\newfontfamily\itfont{Exo2-LightItalic.ttf}
\newfontfamily\bfont{Exo2-Bold.ttf}
\newfontfamily\bifont{Exo2-BoldItalic.ttf}
\newfontfamily\mainf{Exo2-Light.ttf}
\newfontfamily\grfont{FiraSans-LightItalic.ttf}

\definecolor{captcol}{HTML}{495AA8}
\definecolor{ycol}{HTML}{EDAB20}


% Thanks to Douglas Mencken
% See https://tex.stackexchange.com/questions/10001/can-one-draw-a-hyperbola-with-arguments-in-tikz
%
% #1 optional parameters for \draw
% #2 angle of rotation in degrees
% #3 offset of center as (pointx, pointy) or (name-o-coordinate)
% #4 length of plus (semi)axis, that is axis which hyperbola crosses
% #5 length of minus (semi)axis
% #6 how much of hyperbola to draw in degrees, with 90 you’d reach infinity
%
\newcommand\tikzhyperbola[6][thick]{%
    \draw [#1, rotate around={#2: (0, 0)}, shift=#3]
        plot [variable = \t, samples=1000, domain=-#6:#6] ({#4 / cos( \t )}, {#5 * tan( \t )});
    \draw [#1, rotate around={#2: (0, 0)}, shift=#3]
        plot [variable = \t, samples=1000, domain=-#6:#6] ({-#4 / cos( \t )}, {#5 * tan( \t )});
}

\theoremstyle{plain}

\DeclareMathOperator{\N}{\mathbb{N}}
\DeclareMathOperator{\Q}{\mathbb{Q}}
\DeclareMathOperator{\R}{\mathbb{R}}

\DeclareMathOperator{\cl}{cl}
\DeclareMathOperator{\inter}{int}

\DeclareSymbolFont{cyrillic}{T2A}{cmr}{m}{it}
\DeclareMathSymbol{\Sha}{\mathalpha}{cyrillic}{216}
\DeclareMathSymbol{\yy}{\mathalpha}{cyrillic}{251}

\usepackage{xpatch}

\xpatchcmd\swappedhead{~}{.~}{}{}
\swapnumbers

\newtheorem{note}[equation]{Uwaga}
\newtheorem{lemma}[equation]{Lemat}
\newtheorem{definition}[equation]{Definicja}

\setlength{\linewidth}{146mm}
\setlength{\columnsep}{7mm}

\newenvironment{innerproof}
               {\renewcommand{\qedsymbol}{$\blacksquare$}\renewcommand\proofname{Dowód}\proof}
               {\endproof}

               \AtEndPreamble{%
                 \apptocmd{\normalsize}{\fontsize{10pt}{11pt}\selectfont}{}{}
               }%

\addto\captionsbelarusian{\renewcommand{\chaptername}{Разьдзел }}

\begin{document}

\chapter{Уводзіны}

\section{Уступ}

Гэта кніга пра лікі, геамэтрыю і іх сувязь з космасам (і сэнсам жыцьця,
42). Перад тым, як пераходзіць да асноўных тэм, трэба разгледзець некалькі
падставовых матэматычных паняцьцяў, абысьціся безь якіх немагчыма (альбо вельмі
складана) і якія сустракаюцца на кожным кроку ў жыцьці чалавека, зацікаўленага
фізыкай і матэматыкай.

\section{Вытворная}

Пачнем нашую гісторыю з хуткасьці. Як вызначыць хуткасьць машыны? Слова
``вызначыць'' можа мець два значэньні: даць азначэньне і памераць. Фізыкаў
цікавіць і тое, і тое. Перад тым, як нешта вымяраць на практыцы, някепска
паспрабаваць зразумець, што гэта. Зь іншага боку ці мае сэнс гаварыць пра нешта,
што на практыцы памераць немагчыма?

У гэтай кнізе нас цікавіць у першую чаргу матэматыка, таму зоймемся першым
значэньнем. Што такое ``сярэдняя хуткасьць'', сказаць даволі проста: гэта
пройдзеная адлегласьць, падзеленая на час, які машына была ў руху. ``Імгненная
хуткасьць'' --- рэч складанейшая і ёй мы і зоймемся.

Каб спрасьціць сабе жыцьцё, разгледзім машыну, якая рухаецца па бясконцай
простай шашы (для нашых мэтаў хопіць проста вельмі доўгай шашы, каб машына
ніколі не даяжджала да яе канца). Адзначым, што хуткасьць --- рэч адносная і мы
вызначаем яе ў нейкай сыстэме адліку, у нашым выпадку будзе вызначаць яе адносна
шашы.

Выберам пэўны пункт шашы за пачатак і ўвядзем на шашы каардынаты (напрыклад, у
кілямэтрах): адлегласьць улева ад пачатку будзем лічыць адмоўнай, а ўправа ---
дадатнай. Такім чынам, маем функцыю $s(t)$ --- каардыната машыны ў момант часу
$t$.

Рагледзім два моманты часу, $t_1$ і $t_2$ ($t_2 > t_1$), сярэдняя хуткасьць
машыны $v_{s}$ паміж гэтымі момантамі задаецца формулай
\begin{equation*}
  v_{s} = \frac{s(t_2) - s(t_1)}{t_2 - t_1}.
\end{equation*}

Калі нас цікавіць хуткасьць машыны ў вызначаны момант часу $t$, то яе можна
знайсьці набліжана. Іменна, возьмем моманты $t$ і $t + \Delta t$, дзе $\Delta t$
--- маленькі прамежак часу, на працягу якога машына не пасьпявае моцна
разагнацца ці затармазіць. Тады сярэдняя хуткасьць за гэты прамежак часу будзе
прыблізна роўная імгненнай хуткасьці $v_{t}$ ў момант $t$:
\begin{equation*}
  v_{t} \approx \frac{s(t+\Delta t) - s(t)}{\Delta t}.
\end{equation*}

Для машыны, якая стаіць на месцы ў пункце $a$ (дзеля скарачэньня, мы называем
пункты іхнымі каардынатамі), маем $s(t) = a$ для любога моманту часу $t$, так
што
\begin{equation*}
  v_{t} \approx \frac{a - a}{\Delta t} = 0,
\end{equation*}
вынік не залежыць ад часу, таму мы разумеем, што можам узяць у гэтай формуле як
заўгодна маленькі прамежак часу $\Delta t$. Такім чынам, імгненная хуткасьць
машыны, якая стаіць на месцы, дакладна роўная нулю. Гэта вельмі добра, было б
дзіўна, калі б у такой сытуацыі мы атрымалі ненулявую хуткасьць!

Возьмем цяпер машыну, якая рухаецца з сталай хуткасьцю, роўнай $a$, пачаўшы з
пункту $b$ у момант часу $t_0=0$. За час $t$ машына пройдзе $at$ кілямэтраў, так
што яе каардыната ў гэты момант будзе роўная $at + b$. Падставім у нашую формулу
\begin{equation*}
  v_{t} \approx \frac{a(t+\Delta t) + b - at - b}{\Delta t} = a.
\end{equation*}
Вынік зноў не залежыць ад велічыні прамежку часу $\Delta t$, так што і тут
імгенная хуткасьць машыны вылічаная дакладна і, як мы і чакалі, роўная $a$.

Ідзем далей. Дапусьцім, наш кіроўца любіць хуткасьць, а рухавік у яго вельмі
добра, таму ён едзе з сталым паскарэньнем $a$, пачаўшы ў момант часу $t_0=0$ з
пункту $0$ і нулявой хуткасьці. Атрымліваем каардынату ў момант часу $t$ роўную
$at^2$ (чытач павінен ведаць гэта з школьнага курсу фізыкі, з разьдзелу,
прысьвечанаму руху з паскарэньнем), так што нашая формула дае:
\begin{equation*}
  v_{t} \approx \frac{a(t+\Delta t)^2 - at^2}{\Delta t} = \frac{2at\Delta t + (\Delta t)^2}{\Delta t} = 2at + \Delta t.
\end{equation*}
Цяпер сярэдняя хуткасьць залежыць ад прамежку часу $\Delta t$, але заўважым, што
яе даюць два складнікі, адзін незалежны ад $\Delta t$, а другі роўны $\Delta
t$. Другі складнік становіцца ўсё меншым і меншым па меры зьмяншэньня $\Delta
t$, так што вынік набліжаецца да $at$. Так што дакладная імгненна хуткасьць у
момант $t$ будзе роўная папросту $2at$.

Віншую чытача, мы палічылі трэці ліміт у гэтай кнізе! (Першыя два былі трывіяльныя.)

Уважліва паглядзім на тое, што адбылося: мы мералі сярэднюю хуткасьць на ўсё
меншым прамежках часу і заўважылі, што яна набліжаецца да пэўнага ліку. Гэта лік
і будзе імгненнай хуткасьцю.

Ці можна палічыць хуткасьць для больш складанага руху машыны? Альбо, больш
агульна, калі ёсьць нейкі працэс, які адбываецца ў часе і апісваецца нейкай
велічынёй, ці можна палічыць імгненную хуткасьць зьмяненьня гэтай велічыні?
Альбо яшчэ больш агульна, калі ёсьць залежнасьць паміж дзьвюма велічынямі, ці
можна палічыць імгненную хуткасьць зьмяненьня адной велічыні адносна зьмяненьня
іншай?

Удакладнім задачу: ёсьць рэчаісная функцыя $f$ (то бок функцыя, значэньні якой
--- рэчаісныя лікі) зададзеная на рэчаісных ліках, так што для ліку $x$ яе
значэньне роўнае $f(x)$, як знайсьці хуткасьць зьмяненьня гэтай функцыі пры
зьмяненьні $x$?

Сярэднюю хуткасьць зьмяненьня $f(x)$ нам дае формула, прыведзеная
вышэй. Менавіта, возьмем лік $x$ і невялікі лік $\Delta x$, тады сярэдняя
хуткасьць зьмяненьня $f(x)$ у пункце рэчаіснай простай $x$ будзе дадзеная
формулай:
\begin{equation*}
  \frac{f(x + \Delta x)}{\Delta x}.
\end{equation*}

Каб знайсьці імгненную хуткасьць, трэба браць усё меншыя $\Delta x$ і глядзець,
да чаго імкнецца прыведзены вышэй выраз. Каб даведацца гэта, возьмем бясконца
маленькае зьмяненьне аргумэнту $dx$. Але што значыць ``бясконца маленькі''? Ці
мае гэты выраз нейкі сэнс?

Існуюць два адказы на пытаньне пра сэнс выразу ``бясконца маленькі''.

Першы адказ --- ін\-жы\-нэр\-на-фі\-зыч\-ны. Калі мы разглядаем нейкую фізычную
ці практычную задачу, то ў ёй будуць пэўныя характэрныя памеры, у межах якіх
хуткасьці зьмяненьня разгляданых функцыі самі амаль не зьмяняюцца. Калі ўзяць
$\Delta x$ меншым за такі характэрны памер, то $\frac{f(x + \Delta x)}{\Delta
  x}$ будзе набліжаць шуканую хуткасьць вельмі добра, так што з практычнага
пункту гледжаньня $\Delta x$ будзе бясконца маленькім.

Як гэта выглядае ў канкрэтных выпадках? Для машыны на шашы, якая разганяецца з
сталым паскарэньнем, рэалістычнае паскарэньне --- крыху меней за
$3$~м/с$^2$. Такім чынам, за $0,01$~с хуткасьць зьменіцца на $0,03$~м/с, што
складае каля $0,1$~км/г, велічыня для практычных мэтаў блізкая да нуля. Так што
для машыны ў нармальных умовах $0,01$~с можна лічыць бясконца малым часам. А
вось у выпадку астранамічных падзей бясконца малой велічынёй можа аказацца,
напрыклад, адзін сьветлавы год (калі мы пра адлегласьць) ці, адпаведна, адзін
год (калі мы пра час).

Існуе і матэматычны адказ. Тэрміну ``бясконца маленькі'' можна прыпісаць строгае
матэматычнае значэньне. Папярэджваючы здагадкі прасунутых школьнікаў і
экспэртаў, адзначу, што я не пра мэтафарычнае значэньне велічыні, якая імкнецца
да нуля (а імкненьне выражанае, напрыклад, на мове $\epsilon-\delta$), а пра
актуальна бясконца маленькую велічыню. Адпаведная матэматыка была створаная ў
60-х гадах XX стагодзьдзя. Мы пакуль будзем разважаць нефармальна, арыентуючыся
на ўспрыняцьце фізыкаў і інжынэраў, але важна ведаць, што спрашчаючы, мы не
падманваем: матэматыкі могуць сфармуляваць усе нашыя сьцьверджаньні строга.

Карыстаючыся паняцьцем бясконца маленькага ліку, мы можнам даць (амаль)
дакладнае азначэньне імгненнай хуткасьці. Хуткасьць зьмяненьня функцыі $f$ у
пункце $x$ (якую мы пазначаем $f'(x)$) вызначаецца як
\begin{equation*}
  f'(x) = \frac{f(x+dx) - f(x)}{dx},
\end{equation*}
дзе $dx$ --- бясконца маленькая велічыня. Такая хуткасьць называецца вытворнай
функцыі $f$ у пункце $x$.

Велічыня $f(x+dx)-f(x)$ звычайна пазначаецца $df$, а сама вытворная таксама
запісваецца проста як $\frac{df}{dx}$.

Якія ўласьцівасьці маюць бясконца маленькія і як зь іхнай дапамогай знаходзіць
вытворныя? Пачнем з таго, што нас цікавіць канечнае значэньне вытворнай, таму
ўсе бясконца маленькія складнікі ў канчатковым адказе мы можам
ігнараваць. Праілюструем гэта вылічэньнем вытворнай функцыі $f(x) = x^2$ (якое
паўтарае, зробленае вышэй):
\begin{equation*}
  f'(x) = \frac{(x+dx)^2 - x^2}{dx} = \frac{2xdx + (dx)^2}{dx} = 2x + dx = 2x,
\end{equation*}
у канцы мы адкінулі $dx$, бо нас цікавіць канечнае значэньне. Можна ўявіць сабе,
што кожны рэчаісны лік аточаны хмарай бясконца маленькіх, так што, калі
$\epsilon$ --- бясконца маленькі лік, то да хмары, якая атачае, напрыклад, $1$
належаць $1 + \epsilon$, $1 + \epsilon^2$ і г.д. Нас цікавяць канечныя адказы,
таму ў канчатковым выніку мы не адрозьніваем лікаў з гэтай хмары, усе яны для
нас роўныя $1$.

Паколькі ў формуле вытворнай мы дзелім на $dx$, а ў адказе адкідваем бясконца
маленькія складнікі, то ў працэсе вылічэньня мы можам адразу ж адкідваць
складнікі, прапарцыйныя $dx^2$ і вышэйшым ступеням.

Узброіўшыся гэтым правілам, мы можам знайсьці вытворную функцыі $f(x) = x^n$,
дзе $n$ --- натуральны лік:
\begin{equation*}
  f'(x) = \frac{(x+dx)^n - x^n}{dx} = \frac{x^n + nx^{n-1}dx - x^n}{dx} = nx^{n-1},
\end{equation*}
дзе ў лічніку мы скарысталіся пачаткам формулы бінома Ньютана, адкінуўшы ўсе
складнікі з квадратамі і вышэйшымі ступенямі $dx$.

\section{Формулы з вытворнымі}

Цяпер мы можам палічыць вытворную сумы дзьвюх функцый. Калі $h(x) = f(x) +
g(x)$, маем
\begin{multline*}
  h'(x) = \frac{f(x+dx) + g(x+dx) - f(x) - g(x)}{dx} = \\ = \frac{f(x+dx) - f(x)}{dx}
  + \frac{g(x+dx) - g(x)}{dx} = f'(x) + g'(x),
\end{multline*}
так што вытворная сумы роўная суме вытворных, скарочана (не запісваючы аргумэнты
функцый)
\begin{equation*}
  (f + g)' = f' + g'.
\end{equation*}

Што будзе з здабыткам? Для яго атрымаецца складанейшая формула:
\begin{multline*}
  (fg)'(x) = \frac{f(x+dx)g(x+dx) - f(x)g(x)}{dx} = \\ = \frac{f(x+dx)g(x+dx) -
    f(x)g(x+dx) + f(x)g(x+dx)- f(x)g(x)}{dx} = \\ = \frac{(f(x+dx) -
    f(x))g(x+dx)}{dx} + \frac{f(x)(g(x+dx) - g(x))}{dx} = \\ = f'(x)g(x+dx) +
  f(x)g'(x) = f'(x)g(x) + f(x)g'(x).
\end{multline*}\

Тут трэба зьвярнуць увагу на некалькі рэчаў. Па-першае, пазначэньні: $(fg)'(x)$
пазначае значэньне вытворнай здабытку функцый $f$ і $g$ у пункце $x$. Па-другое,
мы скарысталіся распаўсюджаным у матэматыцы прыёмам ``адняць і дадаць'', адняўшы
і дадаўшы $f(x)g(x+dx)$, так, што вынік не зьмяніўся. Нарэшце, у канчатковым
выніку мы праігнаравалі $dx$ у аргумэнце $g(x+dx)$, так рабіць можна, таму што
мы ўжо падзялілі на $dx$ і нас цікавіць канечны вынік, так што мы не
адрозьніваем $x+dx$ і $x$: як бы хутка не зьмянялася функцыя $g$, розьніца паміж
$g(x+dx)$ і $g(x)$ --- бясконца маленькая і можа даць канечны лік толькі пры
дзяленьні на іншы бясконца маленькі лік.

Як вытворная паводзіць сябе пры кампазыцыі функцый? Лічым (кампазыцыю матэматыкі
пазначаюць $f \circ g(x) = f(g(x))$, $f \circ g$ чытаецца ``$f$ пасьля
$g$''):
\begin{equation*}
  (f\circ g)'(x) = \frac{f(g(x+dx)) - f(g(x))}{dx}.
\end{equation*}

Што зрабіць далей? Тут трэба крыху падумаць над тым, чаму роўнае
$f(x+dx)$. Уважліва глядзім на формулу
\begin{equation*}
  f'(x) = \frac{f(x+dx) - f(x)}{dx}
\end{equation*}
і бачым, што
\begin{equation*}
  f(x+dx) = f(x) + f'(x)dx.
\end{equation*}

Чым гэта нам дапаможа? Сачыце за рукамі, карыстаючыся $g(x+dx) = g(x) +
g'(x)dx$, запісваем
\begin{equation*}
  \frac{f(g(x+dx)) - f(g(x))}{dx} = \frac{f(g(x) + g'(x)dx) - f(g(x))}{dx},
\end{equation*}
але ж $g'(x)dx$ такая ж бясконца маленькая велічыня, як і $dx$, так што
\begin{equation*}
  f(g(x) + g'(x)dx) = f(g(x)) + f'(g(x))g'(x)dx.
\end{equation*}

Калі сабраць усё разам, атрымаем
\begin{equation*}
  (f \circ g)'(x) = \frac{f(g(x)) + f'(g(x))g'(x)dx - f(g(x))}{dx} =
  f'(g(x))g'(x).
\end{equation*}

Перамогшы кампазыцыю, можам пазмагацца зь дзяленьнем:
\begin{equation*}
  \left(\frac{1}{f}\right)'(x) = \frac{\frac{1}{f(x+dx)} - \frac{1}{f(x)}}{dx}.
\end{equation*}

\end{document}
